\subsection{Traffic Prediction}
\label{sec:traffic-prediction}

The host CPU prepares circuit candidates before a communication phase
and checks them when its current plan becomes visible.
For Dispatch, we adapt the adjacent-layer expert-load predictor used
by MixNet~\cite{liao2025mixnet}. Let $\mathbf{x}_{\ell-1,s}$ denote the
normalized expert-slot distribution of the preceding layer for tokens
whose current source NI $s$ is already known. Their predicted distribution
is $\widehat{\mathbf{y}}_{\ell,s}=P_\ell\mathbf{x}_{\ell-1,s}$.
The layer-specific matrix is fitted on development requests as
\begin{equation}
  \min_{P_\ell\geq0}\sum_i w_i
  \|\mathbf{y}_i-P_\ell\mathbf{x}_i\|_2^2,
  \qquad \mathbf{1}^{\mathsf T}P_\ell=\mathbf{1}^{\mathsf T},
  \label{eq:noc-load-prediction}
\end{equation}
where $\mathbf{x}_i$ and $\mathbf{y}_i$ describe the same token group in
adjacent layers, and $w_i$ weights its valid slot count and recency.
Matrix $P_\ell$ has $E_\ell$ rows and $E_{\ell-1}$ columns, where
$E_\ell$ is the expert count of layer $\ell$.
Prefill and decode use separate models. A same-layer historical mean
supplies the estimate when preceding-layer information is unavailable.
Known source groups, expert placement, and packing rules convert these
load estimates into source--destination demand; unknown sources remain
unlocated. Once Top-K and packing metadata are visible, actual packet
sizes replace the estimates. Source-read and transmit readiness determine
packet release times.

For Combine, the confirmed Dispatch assignments identify the returning
destinations. Output shape, precision, aggregation, and packing determine
the return packets; expert completion determines their release windows.
Demand history is indexed by layer, phase, prefill/decode, batch range,
and mapping. Measured congestion additionally uses a fingerprint of the
physical topology and effective CS configuration. Completed observations
before the prediction cutoff supply history; a fully observed empty
hotspot set differs from a missing observation. Historical hotspots guide
candidate locations while current demand determines the decision.

\subsection{Reconfiguration Problem}
\label{sec:noc-problem}

\paragraph{Inputs}
For one configuration decision, $\mathcal F$ is the set of directed
source--destination flows and $\mathcal M_f$ the remaining, unadmitted
packets of flow $f$. Packet $m\in\mathcal M=\bigcup_f\mathcal M_f$ carries
$b_m$ useful payload bytes, and $D_f=\sum_{m\in\mathcal M_f}b_m$ is the
flow's remaining volume. The wired topology supplies a set
$\mathcal P_m$ of simple CS paths between that packet's source and
destination; $\mathcal P=\bigcup_m\mathcal P_m$. Each path records its
NI attachments and ordered CS input--output connections. The baseline
Mesh route is denoted by 0. A fixed planning interval of $H$ NoC cycles
and the visible state determine the available service budgets.

\paragraph{Decision variables and objective}
Binary $z_p$ selects CS path $p$. Binary $x_{m,p}$ assigns the whole
packet $m$ to that path. We maximize the useful volume assigned to
the reconfigurable network:
\begin{equation}
  \max_{\mathbf z,\mathbf x}\quad
  V(\mathbf z,\mathbf x)=
  \sum_{m\in\mathcal M}b_m\sum_{p\in\mathcal P_m}x_{m,p}.
  \label{eq:noc-volume-objective}
\end{equation}
Headers consume service but contribute no useful payload to $V$.
Packets without a CS assignment use the Mesh. A flow can therefore
use several CS paths and the Mesh across its packets.

\paragraph{Assignment and switch constraints}
Let $g_m$ indicate that the transfer is eligible for new CS admission:
its source payload is resident, its receive context is allocated, and
it has no existing path binding. The variable domains and assignment
constraints are
\begin{align}
  z_p&\in\{0,1\}, &&p\in\mathcal P,
    \label{eq:noc-binary-variables}\\
  x_{m,p}&\in\{0,1\}, &&m\in\mathcal M,\ p\in\mathcal P_m,
    \nonumber\\
  \sum_{p\in\mathcal P_m}x_{m,p}&\leq1, &&m\in\mathcal M,
    \label{eq:noc-packet-assignment}\\
  x_{m,p}&\leq z_p,\quad x_{m,p}\leq g_m.
    \label{eq:noc-path-activation}
\end{align}
Equation~\eqref{eq:noc-packet-assignment} counts each packet once.
Let $A_{c,p}=1$ when path $p$ traverses CS $c$, and zero otherwise.
The single-connection CS circuit requires
\begin{equation}
  \sum_{p\in\mathcal P}A_{c,p}z_p\leq1,
  \qquad c\in\{0,\ldots,12\}.
  \label{eq:noc-cs-exclusion}
\end{equation}
The physical path table admits wired directions only, excludes a
direction's return to itself and repeated CSs, and terminates at the
specified destination. Source and destination attachments are actual
NI links; nodes never act as intermediate CS forwarding points.

\paragraph{Shared service constraints}
Resources $r\in\mathcal R$ include directed links, NI transmit and
receive service, DMA engines, and SRAM-bank access resources.
The coefficient $u_{m,p,r}$ is packet $m$'s service work on resource $r$
along path $p$, expressed in NoC-cycle equivalents using the resource's
service rate. It includes the actual packet framing; $u_{m,0,r}$ is the
corresponding Mesh cost. Budget $C_r$ is the service available during
$H$ after reserving work for existing bindings, background traffic,
and control messages outside $\mathcal M$. Using the derived Mesh
assignment $x_{m,0}=1-\sum_{p\in\mathcal P_m}x_{m,p}$, the constraint is
\begin{equation}
  \sum_{m\in\mathcal M}\sum_{p\in\mathcal P_m\cup\{0\}}
    u_{m,p,r}x_{m,p}\leq C_r,\qquad r\in\mathcal R.
  \label{eq:noc-resource-capacity}
\end{equation}
The $p=0$ contribution retains the Mesh work of unassigned packets. For a
shared source or destination resource, both route types consume its
service. Different CS attachments may operate concurrently while
sharing these budgets. Queueing, release dependencies, and finite-buffer
backpressure enter the temporal check in the configuration algorithm.

\paragraph{Mapping variables to hardware}
For the wired directions $i\ne o$ of CS $c$, constant $A_{c,i,o,p}$
indicates that path $p$ uses input $i$ and output $o$. The connection
and enable are determined by $\mathbf z$:
\begin{equation}
  h_{c,i,o}(\mathbf z)=\sum_{p\in\mathcal P}A_{c,i,o,p}z_p,
  \qquad e_c=\sum_{i\ne o}h_{c,i,o}(\mathbf z).
  \label{eq:noc-switch-mapping}
\end{equation}
Equation~\eqref{eq:noc-cs-exclusion} leaves at most one active pair.
That pair determines the output code and the input's position in that
output MUX's allowed-input list. When $e_c=0$, the CS is disabled.
The $x_{m,p}$ assignments identify the source-CS attachment for each
packet in the NI transmit queue. These CS states and NI assignments
form configuration $S(\mathbf z,\mathbf x)$, using the register fields
and commands in Section~\ref{sec:noc-config-interface}.

\subsection{Configuration Algorithm}
\label{sec:noc-algorithm}

\paragraph{Path preparation}
Offline, we enumerate local paths through a shared lower CS and
cross-region paths through U-connected lower CSs and either simple
direction around the upper network. Each path stores its direction
codes, source and destination attachments, and resource coefficients.
Before Top-K, predicted demand prepares a tagged candidate cache.
Current packet sizes, eligibility, mapping, and configuration replace
the cached assumptions when the plan becomes visible.

\paragraph{Flow ordering and allocation}
Flows are ordered by decreasing remaining payload $D_f$; a fixed flow
identifier breaks ties. Flows without eligible packets are skipped.
Packets within a flow follow release order,
then transfer identifier and byte offset. For a proposed circuit set,
\textsc{Assign} considers permitted CS paths in increasing unloaded
latency and path-identifier order. It replaces each accepted packet's
Mesh work with its CS work and checks the added service against the
remaining budgets. Packets that cannot use a selected path remain on
the Mesh. An unresolved Mesh budget excess can be reduced during
construction; only complete vectors satisfying
\eqref{eq:noc-binary-variables}--\eqref{eq:noc-resource-capacity}
enter the final candidate pool.

\paragraph{Conflict resolution}
Starting from an empty proposed circuit set, each round considers the
$K_f$ largest remaining flows and at most $K_p$ compatible paths per
flow. These are CPU search budgets. \textsc{Extensions} forms single-path
and mutually compatible two-path additions. Each addition preserves
existing circuit choices and is assigned jointly; two separate single-path
volumes are never added as a substitute for this allocation.
The extension with the largest increase in $V$ becomes the next
construction state. Ties favor fewer changes from the active CS states,
then fixed path identifiers. Construction stops when no volume can be
added, no compatible path remains, or its time budget expires. Because
each successful round reserves a new CS, there are at most 13 rounds.
Thus a two-path extension budget permits configurations containing more
than two circuits. \textsc{Feasible} stores the tested allocations that
satisfy all constraints; the active configuration $S_0$ is also retained.

\paragraph{Reconfiguration check and application}
The feasible pool is checked in decreasing $V$. A tie with the active
configuration retains it; other ties favor fewer state changes.
A time-window model propagates serviced packets
through the selected routes, source reads, receive queues, and destination
bank writes. It includes the candidate's telemetry, CPU computation,
pause, drain, writes, descriptor installation, and resume on a common
timeline. Let $\widehat T(S)$ be the latest estimated completion of the
current wave's required consumers, and $\eta(S)$ a paired-error margin
from independent calibration requests. A proposed update is accepted
only if
\begin{equation}
  \widehat T(S)+\eta(S)<\widehat T(S_0).
  \label{eq:noc-switching-check}
\end{equation}
For background wave $b$, its estimated completion $\widehat T_b$ also
satisfies $\widehat T_b(S)-\widehat T_b(S_0)+\eta_b(S)\leq\epsilon_b$,
where $\eta_b$ is its error margin and $\epsilon_b$ its allowed delay.
These checks control application of the volume-based allocation.
The active network continues during computation. Before applying an
update, the CPU refreshes completed and bound work and rechecks versions,
eligibility, and the timing conditions. If no candidate passes within
the decision budget, it retains $S_0$.

Algorithm~\ref{alg:noc-configuration} gives the steps. The NI protocol
finishes old transfers before changing CS states, installs the packet
assignments with the matching version, and resumes after all required
acknowledgements. Completed observations update demand and
configuration-specific congestion histories separately.

\begin{algorithm}[!t]
  \caption{Volume-based NoC configuration}
  \label{alg:noc-configuration}
  \small
  \begin{algorithmic}[1]
    \Require Current packets $\mathcal M$, budgets $C_r$, cache, active $S_0$
    \Ensure Retain $S_0$ or apply a checked configuration
    \State Sort flows by $D_f$ and packets by release order
    \State $S\gets\varnothing$; $\mathcal C\gets\{S_0\}$
    \While{construction budget remains}
      \State $\mathcal E\gets\Call{Extensions}{S,K_f,K_p}$
      \Statex \Comment{One or two new CS-disjoint paths}
      \For{$S'\in\mathcal E$}
        \State $(\mathbf z',\mathbf x')\gets\Call{Assign}{S',\mathcal M}$
        \If{$\Call{Feasible}{\mathbf z',\mathbf x'}$}
          \State Add $S'(\mathbf z',\mathbf x')$ to $\mathcal C$
        \EndIf
      \EndFor
      \State Choose the extension with largest added $V$
      \If{no positive increase}
        \State \textbf{break}
      \EndIf
      \State $S\gets$ chosen extension
    \EndWhile
    \State Order $\mathcal C$ by $V$, then fewer CS changes
    \For{$S'\in\mathcal C$ in this order}
      \If{$S'=S_0$ or decision budget expired}
        \State \Return $S_0$
      \EndIf
      \State Refresh state; check completion and background delay
      \Statex \Comment{Control and shared-endpoint costs included}
      \If{versions, eligibility, and timing checks pass}
        \State Pause new admissions and finish old transfers
        \State Collect drain responses; write CS states and confirm
        \State Install packet assignments and version; confirm
        \State Resume admissions; \Return $S'$
      \EndIf
    \EndFor
    \State \Return $S_0$
  \end{algorithmic}
\end{algorithm}
